%
A machine-learning classification approach was used to estimate the activity levels of the participants.
%
Two machine-learning models were adopted: a random forest (RF) and an extreme gradient boosting (XGBoost) classifier \cite{ho1995random, 10.1145/2939672.2939785}.
%
Training and evaluation used short-video data.
RF and XGBoost were implemented using the scikit-learn application programming interface (API) \cite{pedregosa2011scikit,buitinck2013api}.
RF combines multiple decision trees and enables feature-importance evaluation during prediction \cite{ho1995random}.
XGBoost applies gradient boosting for high predictive accuracy \cite{10.1145/2939672.2939785}.
%
For each dataset pair, the element-wise absolute differences of point-cloud features were used as input.
  %
  ML derived from comparing the GL of the two datasets in each pair was used as output.
  %
  \deleted{The ML derived from GL is regarded as the ground truth, and the ML estimated by the model based on the features is regarded as the estimated result.}
  \revised{The ML derived from GL was regarded as the ground truth, and the ML estimated by the model based on the features was regarded as the estimated result.}
%
The dataset was randomly split into two parts: 80\% for training and 20\% for evaluation.
%
Hyperparameter tuning via grid search optimized the F1 score for each model \cite{10.5555/2188385.2188395}.
Predefined parameter sets for the number of estimators and maximum tree depth were explored for both RF and XGBoost.
